\documentclass[11pt,oneside,reqno]{amsart}
\usepackage[utf8]{inputenc}

\usepackage{enumitem}


\usepackage{amsthm}
\usepackage{mathtools}
\usepackage{amsmath}
\usepackage{amssymb}
\usepackage[all]{xy}
\usepackage{tikz}
\usepackage{tikz-cd}
\usepackage{mathtools}
\usepackage{mathrsfs}

\usepackage{color}

\def\rcomment#1{{\color{red}[{\sf #1}]}}
\def\bcomment#1{{\color{blue}[{\sf #1}]}}
\def\gcomment#1{{\color{green}[{\sf #1}]}}

%% blackboard bold

\newcommand\Z{{\mathbb Z}}
\newcommand\Q{{\mathbb Q}}
\newcommand\Qbar{{\overline{\mathbb Q}}}
\newcommand\N{{\mathbb N}}
\newcommand\R{{\mathbb R}}
\newcommand\C{{\mathbb C}}

\DeclareMathOperator{\sgn}{\text{sgn}}
\DeclareMathOperator{\per}{\text{Per}}
\DeclareMathOperator{\preper}{\text{Preper}}
\DeclareMathOperator{\tr}{tr}


%% bold face macros

\newcommand\br{{\mathbf r}}
\newcommand\bt{{\mathbf t}}
\newcommand\bu{\mathbf{u}}
\newcommand\bv{\mathbf{v}}
\newcommand\bw{\mathbf{w}}
\newcommand\bx{\mathbf{x}}
\newcommand\by{\mathbf{y}}
\newcommand\bX{\mathbf{X}}
\newcommand\bY{\mathbf{Y}}
\newcommand\bL{\mathbf{L}}

%%%%%%%%%%%%%%%%% environments %%%%%%%%%%%%%%%%%%

\numberwithin{equation}{section}

\newtheorem{theorem}{Theorem}[section]
\newtheorem{lemma}[theorem]{Lemma}
\newtheorem{prop}[theorem]{Proposition}
\newtheorem{cor}[theorem]{Corollary}

\theoremstyle{definition}
\newtheorem{definition}[theorem]{Definition}
\newtheorem{example}[theorem]{Example}
\newtheorem{problem}{Problem} 
\newtheorem{step}{Step}

\theoremstyle{remark}
\newtheorem{remark}[theorem]{Remark}
\newtheorem{notation}[theorem]{Notation}
\newtheorem{question}[theorem]{Question}
\newtheorem{conjecture}[theorem]{Conjecture}

\title{Lecture Notes: Limits and Continuity}
\author[sdf]{S. Peter Jeong}

\begin{document}
\maketitle
\begin{abstract}
    A brief introductory exposition on limits, continuity, and related theorems.
\end{abstract}

\section{Limit of a Function}

\subsection{Examples of Limits} We begin with an intuitive definition of limits of real-valued functions. 
\begin{definition}
    For a real-valued function $f:D\to\R$ defined on some subset $D \subseteq \R$, we say that \textit{the limit of $f(x)$ as $x$ approaches $a$ is $L$} if the value of $f(x)$ gets arbitrarily close to $L$ as $x$ gets closer and closer to $a$, without necessarily reaching $a$. 
\end{definition}
That is, the existence of a limit depends only on the behavior of the function near the point $a$, regardless of whether $f(a)$ is defined. In particular, the limit $\lim_{x\to a} f(x)$ exists if and only if any sequence of numbers $\{x_n\}$ approaching $a$ within some neighorhood of $a$
produces a sequence of numbers $\{f(x_n)\}$ approaching the same value $L$. Since we can approach $a$ from both sides (either right or left) on the real number line, for the overall limit to exist, 
the \textit{right-hand limit} $\lim_{x\to a^+} f(x)$ and \textit{left-hand limit} $\lim_{x\to a^-} f(x)$ must both exist and equal the same value $L$. 

Here we look at some ``nontrivial'' examples (or counterexamples) of limits of functions. 
\begin{example}
    Let $f(x) = \sin \frac{1}{x}$. If we ever try to estimate the limit $\lim_{x\to0} f(x)$ by either graphing or calculating $f(x)$ for very small values of $x$, we quickly run into an issue; the functions seems to ``fluctuate'' across many different values as we take $x = 0.01, 0.00001, 0.0000001,$ and so on. 
    
    In fact, the function indeed oscillates \textit{infinitely often} on any open interval containing the point $0$. In particular, if we let 
    \[x_1 = \frac{1}{\frac{1}{2}\pi+2n\pi}, \qquad x_2 = \frac{1}{\frac{3}{2}\pi + 2m\pi},\]
    then any open interval around $0$ contains the points $x_1$ and $x_2$ for sufficiently large $n$ and $m$, while $f(x_1) = 1$ and $f(x_2) = -1$, meaning that the function oscillates between these two values (which are indeed upper and lower bounds of $f$). 
\end{example}

\begin{example}
    Let $f(x) = \begin{cases}
        1 & \text{if } x \in \Q,\\
        0 & \text{if } x \in \R\setminus\Q.
    \end{cases}$ 
    
    Since there is always an irrational number arbitrarily close to any real number, the function oscillates between $0$ and $1$ infinitely often on any given interval on $\R$ and the limit $\lim_{x\to a} f(x)$ does not exist for any $a \in \R$.
    This function is also known as the \textit{Dirichlet function}. 
\end{example}

\begin{example}
    Let $f(x) = \begin{cases}
        \frac{1}{q} & \text{if } x =\frac{p}{q} \text{ in lowest terms},\\
        0 & \text{if } x \in \R\setminus\Q.
    \end{cases}$
    
    Here $f$ is commonly known as the \textit{Thomae function} or the \textit{ popcorn function}. Take some $a \in (0,1)$. Then for any sequence of real numbers $\{x_n\}$ approaching $a$, we can split the sequence in two parts: (1) a sequence of irrational numbers $\{z'_n\}$ approaching $a$, and (2) a sequence of rational numbers $\{z_n\}$ approaching $a$.
    The first sequence $\{z'_n\}$ of irrational numbers produces a sequence of function values $\{f(z'_n)\}$ that approaches $0$.
    
    The second sequence $\{z_n\}$ requires a little more thinking. If we write $z_n = \frac{p_n}{q_n}$ in lowest terms, then the function values $f(z'_n)$ form a sequence of form $\{\frac{1}{q_n}\}$. 
    Now the question becomes: can we find a sequence of rational numbers $\{z_n\}$ approaching $a$ such that the denominators $q_n$ are bounded above by some constant $M$? If not, then the function values $f(z_n) = \frac{1}{q_n}$ must approach $0$.
    The answer is no; if we take $M$ to be any positive integer, then there are only finitely many rational numbers in $(0,1)$ with denominators at most $M$, so we cannot find a sequence of rational numbers approaching $a$ with bounded denominators.
    Thus the function values $f(z_n)$ must approach $0$ as well, and we conclude that $\lim_{x\to a} f(x) = 0$ for any $a \in (0,1)$.
\end{example}

\begin{example}
    Now let $f(x) = \frac{\sin\theta}{\theta}$ and consider the limit of $f(\theta)$ as $\theta$ approaches $0$. Geometrically, $\sin\theta$ is the length of the vertical segment of the unit vector at angle 
    $\theta$ from the positive $x$-axis. We see that the triangle formed by the origin, the point on the unit circle at angle $\theta$, and the unit vector $(1,0)$ has area $\frac{1}{2}\sin\theta$, while
    the sector of the unit circle at angle $\theta$ has area $\frac{1}{2}\theta$. Since the triangle is contained in the sector, we have 
    \begin{align*}
        \frac{\sin\theta}{2} \leq \frac{\theta}{2}, \quad \text{for } 0 < \theta < \frac{\pi}{2}.
    \end{align*}
    On the other hand, the triangle formed by the origin, the unit vector $(1,0)$, and the point $(1, \tan\theta)$ has area $\frac{1}{2}$. Since this larger right triangle contains the sector, we have the inequality
    \begin{align*}
        \frac{\theta}{2} \leq \frac{\tan\theta}{2}.
    \end{align*}
    Combining the two inequalities, we have
    \begin{align*}
        \frac{\sin\theta}{2} \leq \frac{\theta}{2} \leq \frac{\tan\theta}{2}.
    \end{align*}
    Multiplying through by $\frac{2}{\sin\theta}$, we have
    \begin{align*}
            1 \leq \frac{\theta}{\sin\theta} \leq \frac{1}{\cos\theta}.
    \end{align*}
    for $\theta \in (0, \pi/2)$. Thus $\lim_{x\to 0^+} \frac{\sin\theta}{\theta} = 1$. 

    To show that $\lim_{x\to0^-}\frac{\sin\theta}{\theta} = 1$, notice that $\sin\theta$ is an odd function. Thus after changing the sign of $\theta$ we obtain
    \begin{align*}
        \lim_{x\to0^-}\frac{\sin\theta}{\theta} &= \lim_{x\to0^+}\frac{\sin(-\theta)}{-\theta}
        \\&= \lim_{x\to0^+}\frac{-\sin\theta}{-\theta} \\&= \lim_{x\to0^+}\frac{\sin\theta}{\theta} \\&= 1.
    \end{align*}
    Thus we have $\lim_{x\to0} \frac{\sin\theta}{\theta} = 1$.
\end{example}

\subsection{The Squeeze Theorem} So far, we have seen examples of limits of non-polynomial or non-rational functions that require a more sophisticated definition of a limit. 
We continue to examine limits of functions that cannot be computed directly through standard algebraic manipulation. 
\begin{theorem}[Squeeze Theorem]
    Let $f,g,h$ be real-valued functions defined on some open interval containing $a$, except possibly at $a$ itself. If $f(x) \leq g(x) \leq h(x)$ for all $x$ in the interval, and if $\lim_{x\to a} f(x) = \lim_{x\to a} h(x) = L$, then $\lim_{x\to a} g(x) = L$ as well.
\end{theorem}
Essentially, we are trapping the tricky function $g$ between two functions $f$ and $h$ whose limits are well-defined and have the same value. In that case, the limit of 
the trapped function $g$ is forced to approach the same value as well.

\begin{example}
    Consider the limit $\lim_{x\to0} x^2 \sin\frac{1}{x}$. We can apply the squeeze theorem by noting that $-1 \leq \sin\frac{1}{x} \leq 1$ for all $x$, so we have
\begin{align*}
    -x^2 \leq x^2 \sin\frac{1}{x} \leq x^2.
\end{align*}
Since $\lim_{x\to0} -x^2 = \lim_{x\to0} x^2 = 0$, we conclude that $\lim_{x\to0} x^2 \sin\frac{1}{x} = 0$ as well.  
\end{example}


\end{document}